% ============================================================
% Section 1: Introduction
% ============================================================

\section{Introduction}
\label{sec:intro}

Identity document fraud poses a growing threat to national security, financial
integrity, and border control worldwide. The widespread availability of
image-editing software has made it possible to produce convincing forgeries of
passports, national identity cards, and driver's licenses at low cost.
Automated forensic analysis of identity documents
is therefore a critical component of modern digital identity verification
pipelines.

Detection of forged identity documents is inherently a two-level problem. At the
\emph{document level}, a system must decide whether a presented document is
genuine or has been tampered with (binary detection). At the \emph{pixel level},
it must also identify and delineate which specific regions have been altered
(forgery localization). Both objectives are operationally necessary.

Binary detection alone is insufficient for real-world forensic deployments.
A positive detection decision provides no \emph{actionable information} to the
investigator: it does not reveal whether the fraudster altered the portrait
photograph, falsified textual fields, or combined multiple manipulations. Without
spatial evidence, a human expert cannot determine the nature of the fraud or
build a legal case. Detection systems are also increasingly subject to
\emph{explainability requirements} under emerging AI regulations: a binary
``tampered'' label cannot satisfy the obligation
to provide a meaningful explanation of the automated decision. Conversely,
localization alone does not yield a reliable document-level confidence score
suitable for automated accept/reject decisions. A practical forensic system must
address both tasks simultaneously.

The primary technical difficulty in joint detection and localization is that the
two objectives are in \emph{structural tension}: the classification head
benefits from global image features, while the segmentation head requires
spatially precise, pixel-level representations. Prior work either treats the
two tasks independently \citep{trufor2023} or collapses both losses into a
scalar objective via a fixed weight $\lambda$~\citep{kendall2018multi}. The choice of $\lambda$ is arbitrary and dataset-dependent.
A second challenge is the risk of over-fitting the model-selection criterion:
hyperparameter optimization (HPO) that maximizes PR-AUC may sacrifice Dice, and vice versa.

This paper introduces DocVerify, a multi-task forensic document
analysis system that addresses both challenges. Our contributions are:

\begin{enumerate}
  \item \textit{Multi-objective Pareto HPO for multitask forensics.}
    Both PR-AUC and Dice are treated as simultaneous optimization objectives.
    The optimal configuration is selected from the Pareto front by minimizing
    the Euclidean distance to the ideal point~$(1,~1)$, requiring no manual
    task-weight tuning.

  \item \textit{DocVerify architecture.}
    An end-to-end multi-task CNN comprising a lightweight convolutional encoder
    coupled to a U-Net decoder for pixel-level mask prediction and a
    classification head for document-level detection, trained jointly with a
    combined BCE+Dice segmentation loss.

  \item \textit{Rigorous nested cross-validation.}
    A 10-fold outer $\times$ 5-fold inner nested cross-validation scheme with
    50 Multiobjective Tree-structured Parzen Estimator (MOTPE) trials per
    outer fold (2\,500 HPO evaluations total). Final performance is estimated
    on a blind holdout using 30 independent seeds, with Wilcoxon signed-rank
    tests~\citep{wilcoxon1945} and Holm--Bonferroni correction~\citep{holm1979}.
\end{enumerate}

The source code and experimental artifacts are publicly available at
\url{https://github.com/HernanDiaz/IDVerify/}.

\begin{figure*}[!t]
  \centering
  \includegraphics[width=\textwidth]{fig0_architecture_compact}
  \caption{DocVerify architecture. The encoder maps a $3{\times}224{\times}224$
    input to a $256{\times}7{\times}7$ bottleneck through six convolutional
    blocks ($8\!\to\!256$ channels). The U-Net decoder reconstructs the forgery
    mask via skip connections; the classification head outputs a document-level
    tampered/genuine score.}
  \label{fig:architecture}
\end{figure*}

The remainder of the paper is organized as follows. Section~\ref{sec:related}
reviews related work. Section~\ref{sec:method} describes the DocVerify
architecture and HPO framework. Section~\ref{sec:experiments} presents
experimental results. Section~\ref{sec:discussion} discusses the findings and
limitations. Section~\ref{sec:conclusion} concludes.
